\documentclass[10pt,a4paper]{amsart}
\usepackage[margin=19mm]{geometry}
\usepackage{amsmath,amssymb,mathtools}
\usepackage[scr=rsfs,cal=euler]{mathalfa}
\usepackage{xfrac}
\usepackage[unicode,hidelinks,bookmarks=false]{hyperref}
\newcommand{\ie}{i.e.\ }
\newcommand{\cf}{cf.\ }
\newcommand{\resp}{respectively\ }
\newcommand{\Subsection}{\subsection}
\providecommand{\nobreakdash}{\nobreak\hbox{-}}
\setlength{\emergencystretch}{2em}
\multlinegap0pt
\usepackage{bm}
\usepackage{bbm}
\usepackage{dsfont}
\usepackage{xspace}
\usepackage{cancel}
\usepackage{mathtools}
\usepackage{amsthm}
\makeatletter
\@ifundefined{lemma}{\newtheorem{lemma}{Lemma}}{}
\@ifundefined{theorem}{\newtheorem{theorem}{Theorem}}{}
\makeatother
\title[On the Natural Gradient of the Evidence Lower Bound]{On the Natural Gradient of the Evidence Lower Bound}
\author{Nihat Ay, Jesse van Oostrum, Adwait Datar}
\date{Source: arXiv:2307.11249v2 (CC BY 4.0)}
\begin{document}
\maketitle

\noindent\emph{Source and licence.} On the Natural Gradient of the Evidence Lower Bound. Version arXiv:2307.11249v2 (CC BY 4.0), distributed under the Creative Commons Attribution 4.0 International licence, as recorded in the arXiv licence metadata for this exact version and shown on the abstract page https://arxiv.org/abs/2307.11249v2. Full rights notice: licenses/arxiv-2307.11249-CC-BY-4.0.txt.\par\smallskip

\noindent\textit{Editorial scope.} Counted unit: the Theorem whose source label is \texttt{mainthdist} in the article (source0.tex lines 763--776, source label \texttt{mainthdist}), delivered together with the article's own complete original proof of it (source0.tex lines 777--809, one \texttt{proof} environment of 33 lines). No mathematics is written, repaired or re-derived: statement and proof are the article's own text line for line. Required context retained uncounted: invcylind (lines 453--462); defproj (lines 519--521); graddeppr (lines 553--574); gradientinp (lines 555--557); invariancedist (lines 569--572). Every definition, display and label of those lines is retained unchanged. Omitted from this package and not redistributed: 1--452, 463--518, 522--552, 558--568, 573--762, 810--3492. Nothing inside a retained excerpt is omitted or altered: every retained body is delivered exactly as the article's lines are written, so no omission receipt of any kind accompanies this package. Citations: the retained excerpts carry no citation command at all, so no bibliography entry of the article is transcribed and no reference is redistributed. Reference adaptation: none was needed and none was made; every reference inside the delivered unit resolves inside it, so no unresolved reference or citation is produced. Printed slips kept literal: no printed word, symbol, hypothesis or number of the counted statement or of its proof was altered. Layout and class: the article's own class is \texttt{article} with packages including jmlr2e, amsmath, xcolor, tikz, pgfplots, mymacros, comment, xy, lastpage; the wrapper is the lane's pinned \texttt{amsart} editorial wrapper at 10pt on A4, which restates the theorem-like environments the retained text needs and adds the article's own macro definitions that the retained text needs (none), with extra wrapper packages (bm, bbm, dsfont, xspace, cancel, mathtools). The delivered numbering is the unit's own. Assets: no retained span contains a figure environment, a graphics-inclusion command, a PDF-page inclusion, a table or a TikZ picture; no image file is copied into this package, the package declares no asset dependency and no image file is redistributed with it. Licence: version arXiv:2307.11249v2 of this article is distributed under the Creative Commons Attribution 4.0 International licence, as recorded in the arXiv licence metadata for this exact version and shown on the abstract page https://arxiv.org/abs/2307.11249v2. A full rights notice is delivered at licenses/arxiv-2307.11249-CC-BY-4.0.txt. Characters: every retained excerpt is pure ASCII, so the delivered engine, XeLaTeX with its default Latin Modern text font, reports no missing character.
\begin{lemma} \label{invcylind}
Let ${\mathcal M}$ be a cylindrical model in ${\mathcal P}$, let $p$ be a non-singular point of ${\mathcal M}$, and let $\Pi_p$ denote the orthogonal projection of $T_p {\mathcal P}$ onto $T_p {\mathcal M}$. 
Then,  
\begin{equation*}
   \Pi_p ({\mathcal H}_p) 
   \; \subseteq \; {\mathcal H}_p, \qquad 
    \Pi_p ({\mathcal V}_p) 
    \; \subseteq \; {\mathcal V}_p
\end{equation*}
\end{lemma}

\begin{equation} \label{defproj}
   \pi_{\mathcal Q}(p)(x_V, x_H)  \; = \; p^\ast(x_V) p(x_H | x_V).  
\end{equation}

\begin{theorem} \label{graddeppr}
{\bf (a)} Consider first the function $D({\mathcal Q} \| \cdot )$ on ${\mathcal P}$. Then  
\begin{equation} \label{gradientinp}
     {\rm grad}^{\mathcal P}_p  D({\mathcal Q} \| \cdot ) \; = \; p - \pi_{\mathcal Q}(p), 
\end{equation}
where $\pi_{\mathcal Q}(p)$ is defined by (\ref{defproj}).
In order to obtain the gradient of 
$D({\mathcal Q} \| \cdot )$ in a non-singular point $p \in {\mathcal M}$, we have to project (\ref{gradientinp}) onto $T_p {\mathcal M}$, that is 
\begin{equation} \label{gradientOnModel}
     {\rm grad}^{\mathcal M}_p  D({\mathcal Q} \| \cdot )  \; = \; \Pi_p (p - \pi_{\mathcal Q}(p)), 
\end{equation}
where $\Pi_p$ denotes the orthogonal projection $T_p{\mathcal P} \to T_p {\mathcal M}$ with respect to the Fisher-Rao metric on ${\mathcal P}$.
\smallskip

\noindent
{\bf (b)} If ${\mathcal M}$ is cylindrical and $p \in {\mathcal M}$ admissible then 
\begin{equation} \label{invariancedist}
  d\pi_V \left( {\rm grad}^{\mathcal M}_p  D({\mathcal Q} \| \cdot ) \right) 
  \; = \; {\rm grad}^{{\mathcal M}_V}_{\pi_V(p)}  D(p^\ast \| \cdot ).
\end{equation} 
In particular, the equality (\ref{invariancedist}) holds in all points of the maximal model ${\mathcal M} = {\mathcal P}$ where ${\mathcal M}_V = {\mathcal P}_V$. 
\end{theorem}

\begin{theorem} \label{mainthdist}
Let ${\mathcal M}$ be a cylindrical model in ${\mathcal P}$, let $p \in {\mathcal M}$ be admissible, and let $q \in {\mathcal Q}$. Then
\begin{equation}  
 d\pi_V \left( {\rm grad}^{{\mathcal M}}_p \, 
   {\rm GAP}(q, \cdot ) \right) 
   \; = \; 0, \label{vanishgap} 
\end{equation}
and therefore 
\begin{equation} 
   d\pi_V \left( {\rm grad}^{{\mathcal M}}_p 
   D(q \| \cdot ) \right) 
           \; = \;  {\rm grad}_{\pi_V(p)}^{{\mathcal M}_V} \, D(p^\ast \| \cdot). \label{extension}
\end{equation}  
\end{theorem} 

\begin{proof} 
We begin with the gradient of ${\rm GAP}(q,\cdot)$:
\begin{equation*}
  {\rm grad}^{{\mathcal M}}_p \, {\rm GAP}(q, \cdot ) 
        \; = \; \Pi_p \left( {\rm grad}^{{\mathcal P}}_p \,
        {\rm GAP}(q,\cdot ) \right) 
        \; = \; \Pi_p \left( 
        \pi_{\mathcal Q}(p) - q
        \right) 
        \; = \;  \Pi_p \left( {(p - q)}^{\mathcal V} \right). 
\end{equation*}
We know, by definition, that 
${(p - q)}^{\mathcal V}$ is contained in ${\mathcal V}_p$. 
According to Lemma \ref{invcylind}, the vector ${(p - q)}^{\mathcal V}$ remains in  ${\mathcal V}_p$ after projecting it onto the tangent space $T_p {\mathcal M}$, that is,  
\[
   \Pi_p \left( {(p - q)}^{\mathcal V} \right) \; \in \; {\mathcal V}_p. 
\]
Therefore, this vector is mapped via $d\pi_V$ to $0$, which proves 
(\ref{vanishgap}). 
With this, we have   
\begin{eqnarray*}
  d\pi_V \left( {\rm grad}^{{\mathcal M}}_p \, D(q \| \cdot ) \right)  
        & = & 
   d\pi_V \left( {\rm grad}^{{\mathcal M}}_p \, 
   D({\mathcal Q} \| \cdot ) \right) +   
   d\pi_V \left( {\rm grad}^{{\mathcal M}}_p \, 
   {\rm GAP}(q, \cdot )\right) \\
   & = & d\pi_V \left( {\rm grad}^{{\mathcal M}}_p \, 
   D({\mathcal Q} \| \cdot ) \right), 
\end{eqnarray*}
and with (\ref{invariancedist}) of Theorem  \ref{graddeppr} 
we finally obtain (\ref{extension}).
\end{proof}

\end{document}
